
% !TeX root = main.tex
\newpage
\clearpage
\begin{center}
\Large\textbf{Conventions Section\\ (included only in rough draft. For documenting standards, outline, and rules)}
\end{center}

\section*{Document Review Approach}

The objective of the review process is to improve clarity, strengthen
the narrative structure, and reduce unnecessary complexity while
preserving all information required to support the paper's central
argument.

Each paragraph is evaluated against the current manuscript outline to
ensure that it contributes directly to the progression

[
\text{Problem}
$\rightarrow$
\text{Limitation}
$\rightarrow$
\text{Architectural Contribution}
$\rightarrow$
\text{Validation}
$\rightarrow$
\text{Capability Enabled}.
]

For every paragraph, the following questions are applied:

\begin{enumerate}

\item Does the paragraph belong in the section defined by the outline?

\item What information does the paragraph contribute to the narrative?

\item Can the paragraph be shortened without removing information?

\item Can the language be made more objective?

\item Can statements describing quality or significance be replaced by
descriptions of what was done?

\item Does the final sentence naturally motivate the first sentence of
the following paragraph?

\item Does the paragraph introduce new information, or does it repeat
information already presented elsewhere?

\end{enumerate}

Paragraphs are then classified as:

\begin{itemize}

\item \textbf{KEEP} -- The paragraph supports the outline and requires
only minor editing.

\item \textbf{REWRITE} -- The paragraph contains useful information but
requires restructuring, shortening, or improved transitions.

\item \textbf{MERGE} -- The paragraph overlaps substantially with
adjacent material and should be combined with another paragraph.

\item \textbf{MOVE} -- The paragraph belongs in a different section of
the manuscript.

\item \textbf{DELETE} -- The paragraph does not contribute information
required by the outline or repeats information already presented.

\end{itemize}

Particular attention is given to the following editorial principles:

\begin{itemize}

\item Can this sentence be shorter?

\item Can it be more objective?

\item Can it state what was done rather than how good it is?

\end{itemize}

The goal is to produce concise, technically precise prose that allows
the experimental results and analysis to establish significance rather
than relying on subjective claims.


\section{Manuscript Outline}

\begin{verbatim}
Abstract
What was done?
Why is it important?
What was achieved?

1. Introduction
	Why are particle systems important?
		Paragraphs 1–2

	What are particle systems used for?
		Paragraphs 1–2

	Why are large particle systems computationally limited?
		Paragraphs 3–5

	What is the contribution of gPCD?
		Paragraphs 6–7

	Paragraph 1

	   Purpose:
	   - Define particle methods.
	   - Explain why particle methods are useful.
	   - Introduce emergent macroscopic behavior.

	   Contains:
	   - Particle methods model discrete interacting entities.
	   - Applicable to solids, fluids, granular materials,
		 multiphase flows, molecular systems, and other applications.
	   - Individual entities are represented explicitly.
	   - Macroscopic behavior emerges from local interactions.
	   - Particle methods are well suited for studying collective behavior.

	   Presentation Bullets:
	   - Discrete interacting particles.
	   - Local interactions create emergent behavior.
	   - Applicable to many scientific and engineering systems.

	Paragraph 2

	   Purpose:
	   - Explain why large particle simulations are important.
	   - Motivate particle simulation as a discovery tool.
	   - Connect simulation to scientific advancement and AI.

	   Contains:
	   - Particle simulations support scientific discovery.
	   - New physical behavior may only emerge in large-scale simulations.
	   - Machine-learning methods require representative training data.
	   - Simulations generate training datasets.
	   - Simulations enable investigation of poorly understood phenomena.
	   - Simulations complement experiments and analysis.

	   Presentation Bullets:
	   - Large-scale simulations enable scientific discovery.
	   - Emergent phenomena may require many interacting particles.
	   - Simulations generate datasets for machine learning.
	   - Simulations help investigate previously unobserved behavior.

	Paragraph 3

	   Purpose:
	   - Explain why large particle counts are desirable.
	   - Explain why large particle simulations become expensive.
	   - Introduce the need for efficient particle-processing methods.

	   Contains:
	   - Simulation fidelity increases with particle count.
	   - Realistic simulations may require millions or billions of particles.
	   - Larger particle counts improve spatial resolution.
	   - Larger particle counts improve geometric representation.
	   - Larger particle counts enable multiscale investigations.
	   - Increased particle counts substantially increase computational cost.
	   - Efficient particle-processing algorithms are required for large-scale simulation.

	   Presentation Bullets:
	   - High-fidelity simulations require millions to billions of particles.
	   - Larger particle counts improve resolution and physical realism.
	   - Computational cost grows rapidly with problem size.
	   - Efficient particle-processing algorithms are essential.
   
	Paragraph 4

		Purpose:
		- Identify collision detection as a major computational bottleneck.
		- Explain why naïve collision detection is impractical.
		- Introduce broad-phase collision detection.
		- Motivate the need for efficient parallel broad-phase methods.

		Contains:
		- Collision detection is a primary limitation of large-scale particle simulation.
		- Naïve collision detection scales as O(N²).
		- Direct all-to-all evaluation is impractical for large particle counts.
		- Broad-phase methods reduce collision-detection cost.
		- Broad-phase methods restrict testing to nearby particles.
		- Efficient broad-phase construction remains challenging on massively parallel hardware.

		Presentation Bullets:
		- Collision detection is a major performance bottleneck.
		- Naïve collision detection scales as O(N²).
		- Broad-phase methods reduce candidate pairs.
		- Efficient parallel broad-phase construction remains challenging.
	   
	Paragraph 5

		Purpose:
		- Introduce GPUs as a platform for large-scale particle simulation.
		- Explain that hardware capability alone does not solve the problem.
		- Identify occupancy generation as the remaining bottleneck.
		- Transition toward the architectural motivation for gPCD.

		Contains:
		- GPUs provide high computational throughput.
		- GPUs provide high memory bandwidth.
		- GPUs are well suited for particle simulation.
		- Spatial occupancy structures remain difficult to construct efficiently.
		- Occupancy generation remains a bottleneck at large scales.

		Presentation Bullets:
		- GPUs provide massive parallelism.
		- GPUs offer high memory bandwidth.
		- Particle simulation maps naturally to GPU hardware.
		- Occupancy generation remains a major bottleneck.  
   
	Paragraph 6

		Purpose:
		- Introduce gPCD.
		- State the primary contribution of the work.
		- Summarize the key architectural ideas.
		- Explain the significance of the framework.

		Contains:
		- gPCD is a GPU-driven particle collision-detection framework.
		- gPCD uses the graphics pipeline for occupancy generation.
		- gPCD constructs a uniform cell-array occupancy structure.
		- Occupancy generation is organized around particle corner locations.
		- GPU rasterization hardware is used for broad-phase construction.
		- The framework maintains a simple and scalable data structure.
		- The framework supports large particle counts on a single CPU–GPU system.
		- The framework provides a foundation for future large-scale particle simulations.

		Presentation Bullets:
		- GPU-driven particle collision detection (gPCD).
		- Graphics-pipeline occupancy generation.
		- Corner-based occupancy construction.
		- Uniform cell-array broad phase.
		- Scalable single CPU–GPU execution.
	
	Paragraph 7

		Purpose:
		- Clearly define the novelty of the work.
		- Distinguish gPCD from existing occupancy structures.
		- Explain the architectural contribution.
		- Explain why the graphics pipeline is used.

		Contains:
		- The contribution is an execution architecture rather than a new data structure.
		- gPCD integrates occupancy generation directly into the graphics pipeline.
		- Rendering and occupancy construction occur concurrently.
		- Occupancy generation produces the potentially colliding set (PCS).
		- Graphics-pipeline resources participate in collision detection.
		- Resources traditionally reserved for rendering are repurposed for occupancy generation.

		Presentation Bullets:
		- Novelty: execution architecture, not a new data structure.
		- Graphics-pipeline occupancy generation.
		- Concurrent rendering and occupancy construction.
		- PCS generation integrated into rendering workflow.
		- Repurposes graphics resources for collision detection.
	
	
	
	
2. Background and Related Work
   What must the reader know first?

   Section Introduction

	   Purpose:
	   - Introduce the scope of the background section.
	   - Explain that collision-detection methods differ in both
		 data structures and execution architectures.
	   - Establish the relationship between collision detection
		 and rendering.
	   - Prepare the reader for the architectural gap discussion.

	   Contains:
	   - Many collision-detection approaches have been proposed.
	   - Approaches differ in both data structures and execution architectures.
	   - Collision detection may be organized differently relative to rendering.
	   - The section reviews methodologies, architectures,
		 and rendering strategies.
	   - The section identifies the architectural gap addressed by gPCD.

	   Presentation Bullets:
	   - Collision-detection methods vary in structure and execution model.
	   - Rendering and collision detection are often separated.
	   - Existing architectures leave an unexploited opportunity.
	   - Background motivates the architectural contribution of gPCD.
   
2.1 Collision-Detection Frameworks
		What is occupancy generation and why is it needed?

	Paragraph 1

	   Purpose:
	   - Introduce the broad-phase/narrow-phase organization.
	   - Define the two major stages of collision detection.
	   - Establish where gPCD contributes.
	   - Explain why broad-phase generation is important.

	   Contains:
	   - Collision detection is commonly divided into broad-phase and narrow-phase stages.
	   - Broad phase generates collision candidates.
	   - Narrow phase evaluates candidate collisions.
	   - This organization is widely used across collision-detection frameworks.
	   - gPCD's primary contribution occurs in broad-phase occupancy generation.
	   - Narrow phase completes the overall collision-detection process.

	   Presentation Bullets:
	   - Collision detection is typically divided into two stages.
	   - Broad phase generates collision candidates.
	   - Narrow phase evaluates candidate collisions.
	   - gPCD focuses on broad-phase occupancy generation.
	   

	Paragraph 2

		Purpose:
		- Review common GPU broad-phase methodologies.
		- Show that many approaches exist for candidate generation.
		- Explain the common objective of broad-phase methods.
		- Transition from collision-detection methods to execution architectures.

		Contains:
		- Numerous GPU broad-phase methods have been proposed.
		- Common approaches include:
		  - Uniform grids
		  - Spatial hashing
		  - Hierarchical spatial subdivision
		  - Sweep and prune
		  - Image-space methods
		  - Other acceleration techniques
		- Methods differ in how spatial information is organized and queried.
		- All broad-phase methods seek to reduce narrow-phase evaluations.
		- Research has expanded from algorithms to execution architectures.
		- GPU computing has increased interest in architectural design.

		Presentation Bullets:
		- Many GPU broad-phase methodologies exist.
		- Methods differ in spatial organization and query strategy.
		- All seek to reduce narrow-phase collision tests.
		- Recent work increasingly focuses on execution architectures.

2.2 Collision-Detection Architectures
		How have rendering and collision detection traditionally
		been organized?

   Paragraph 1

	   Purpose:
	   - Introduce execution architectures as distinct from collision-detection methods.
	   - Describe the historical evolution of collision-detection execution.
	   - Explain the migration from CPU to GPU implementations.
	   - Establish the architectural spectrum considered in the literature.

	   Contains:
	   - Collision-detection methodologies can be implemented using different execution architectures.
	   - Early systems performed broad-phase and narrow-phase operations on the CPU.
	   - Collision detection and rendering were traditionally separate processes.
	   - Increasing particle counts motivated GPU acceleration.
	   - Programmable GPUs enabled portions of the collision-detection pipeline to migrate to the GPU.
	   - GPU parallelism became a major driver of architectural evolution.
	   - Existing architectures span:
		   - CPU-resident implementations
		   - GPU-assisted implementations
		   - Fully GPU-resident implementations

	   Presentation Bullets:
	   - Execution architecture is distinct from collision-detection methodology.
	   - Early systems were CPU-based.
	   - Collision detection and rendering were traditionally independent.
	   - GPU programmability enabled migration of collision detection to the GPU.
	   - Architectures now range from CPU-resident to fully GPU-resident.

	Paragraph 2

		Purpose:
		- Define the major architectural categories.
		- Explain how work is distributed between CPU and GPU.
		- Describe the progression toward increased GPU utilization.
		- Introduce data movement as an architectural concern.

		Contains:
		- CPU-resident architectures perform broad-phase and narrow-phase operations on the CPU.
		- Simulation state is transferred to the rendering system after collision processing.
		- GPU-assisted architectures offload selected collision-detection tasks to the GPU.
		- CPU-managed execution and data structures remain part of the workflow.
		- GPU-resident architectures perform broad-phase and narrow-phase operations entirely on the GPU.
		- GPU-resident approaches reduce host-device communication.
		- GPU-resident approaches improve parallel utilization.
		- GPU-resident approaches reduce data movement overhead.

		Presentation Bullets:
		- CPU-resident: collision detection performed entirely on the CPU.
		- GPU-assisted: selected operations offloaded to the GPU.
		- GPU-resident: broad phase and narrow phase executed on the GPU.
		- Increasing GPU utilization reduces data-transfer overhead.


	Paragraph 3

		Purpose:
		- Identify a common limitation across existing architectures.
		- Explain that rendering and collision detection are usually separate.
		- Highlight the traditional role of the rendering pipeline.
		- Motivate examination of rendering-collision integration.

		Contains:
		- Existing architectures differ in execution location and data movement.
		- Collision detection and rendering are typically separate processes.
		- Collision-detection work is performed in dedicated computational stages.
		- The rendering pipeline primarily consumes simulation state.
		- Rendering generally does not participate in collision detection.
		- This separation motivates investigation of rendering-collision interaction.

		Presentation Bullets:
		- Most architectures separate rendering and collision detection.
		- Collision detection uses dedicated compute stages.
		- Rendering typically consumes simulation results.
		- Rendering rarely contributes to collision detection.
		- This separation motivates the gPCD architecture.   
   
2.3 Rendering and Collision Detection

	Paragraph 1

		Purpose:
		- Define the traditional roles of collision detection and rendering.
		- Explain why the two systems have historically been separated.
		- Establish the conventional simulation-to-rendering workflow.
		- Prepare the reader for alternative rendering-based architectures.

		Contains:
		- Collision detection identifies particle interactions.
		- Collision detection advances simulation state.
		- Rendering visualizes simulation results.
		- Collision detection and rendering perform different functions.
		- Rendering is typically positioned downstream of simulation.
		- Rendering primarily consumes simulation state.
		- Rendering traditionally does not participate in simulation processing.

		Presentation Bullets:
		- Collision detection computes simulation behavior.
		- Rendering visualizes simulation results.
		- Rendering is typically downstream of simulation.
		- Rendering acts primarily as a consumer of simulation state.



	Paragraph 2

		Purpose:
		- Explain how GPU adoption affected rendering and collision detection.
		- Show that GPU execution did not fundamentally change their relationship.
		- Reinforce the separation between rendering and collision detection.
		- Identify the underutilization of rendering resources for collision detection.

		Contains:
		- Rendering and collision detection increasingly leverage GPU parallelism.
		- GPU-resident implementations became common.
		- Rendering and collision detection generally remain logically separated.
		- Collision detection executes in dedicated computational stages.
		- Rendering continues to generate visual output from simulation data.
		- Rendering contributes little or no information to collision detection.
		- GPU adoption alone does not integrate rendering and collision detection.

		Presentation Bullets:
		- Both rendering and collision detection moved to the GPU.
		- GPU residency did not eliminate their separation.
		- Collision detection uses dedicated compute stages.
		- Rendering remains focused on visualization.
		- Rendering contributes little to collision detection.




	Paragraph 3

		Purpose:
		- Review prior attempts to integrate rendering and simulation.
		- Acknowledge graphics-assisted approaches.
		- Explain the remaining limitation of existing methods.
		- Identify the architectural gap that motivates gPCD.

		Contains:
		- Prior work has explored rendering-simulation integration.
		- Approaches include:
			- Image-space techniques
			- GPU data structures
			- Graphics-assisted computation
		- Existing approaches still treat occupancy generation as a dedicated
		  collision-detection task.
		- Candidate construction remains independent of rasterization.
		- Rendering remains primarily downstream of simulation.
		- Rendering rarely acts as an active participant in collision detection.
		- An architectural opportunity remains unexploited.

		Presentation Bullets:
		- Prior work explored graphics-assisted simulation.
		- Image-space and GPU-assisted approaches exist.
		- Occupancy generation remains separate from rasterization.
		- Rendering is still primarily downstream of simulation.
		- An architectural gap remains.



	Paragraph 4

		Purpose:
		- Formulate the central architectural question.
		- Transition from prior work to the gPCD motivation.
		- Introduce graphics-pipeline occupancy generation as a possibility.
		- Establish concurrency between rendering and occupancy generation as a design objective.

		Contains:
		- Existing architectures separate occupancy generation from rendering.
		- This separation suggests an unexplored architectural opportunity.
		- Broad-phase occupancy generation may be integrated into the graphics pipeline.
		- Occupancy construction and rendering may proceed concurrently.
		- Graphics-pipeline resources may contribute directly to collision detection.

		Presentation Bullets:
		- Can occupancy generation be integrated into the graphics pipeline?
		- Can rendering and occupancy construction execute concurrently?
		- Can graphics resources participate directly in collision detection?
		- This question motivates gPCD.


2.4 Gap in Existing Methods
		What architectural capability is missing?

	   Paragraph 1

		   Purpose:
		   - Summarize the findings of the previous background sections.
		   - Restate the traditional relationship between rendering and collision detection.
		   - Identify occupancy generation as a dedicated operation.
		   - Prepare the formal statement of the architectural gap.

		   Contains:
		   - Rendering and collision detection have traditionally remained separate.
		   - This separation persists in GPU-resident implementations.
		   - Rendering primarily consumes simulation state.
		   - Broad-phase occupancy generation is performed separately.
		   - Occupancy generation is treated as a dedicated collision-detection task.
		   - Existing architectures do not integrate occupancy generation into rendering.

		   Presentation Bullets:
		   - Rendering and collision detection remain separate.
		   - This separation persists on modern GPUs.
		   - Rendering consumes simulation state.
		   - Occupancy generation remains a dedicated collision-detection operation.



	Paragraph 2

		Purpose:
		- State the architectural gap explicitly.
		- Define the missing capability in existing methods.
		- Introduce the gPCD solution.
		- Connect the Background section to the Methodology section.

		Contains:
		- Existing methods do not integrate broad-phase occupancy generation directly into the graphics pipeline.
		- This absence constitutes the architectural gap.
		- gPCD addresses the gap through graphics-pipeline occupancy generation.
		- Occupancy construction and rendering proceed concurrently.
		- Graphics-pipeline execution produces the potentially colliding set (PCS).
		- The PCS is used during narrow-phase evaluation.
		- gPCD integrates rendering and occupancy generation into a unified architecture.

		Presentation Bullets:
		- Missing capability:
			- Graphics-pipeline occupancy generation.
		- gPCD integrates occupancy generation directly into the graphics pipeline.
		- Occupancy construction and rendering execute concurrently.
		- Graphics-pipeline execution produces the PCS.
		- gPCD closes the architectural gap.



	Paragraph 3

		Purpose:
		- Clarify the nature of the contribution.
		- Distinguish architectural novelty from algorithmic novelty.
		- Prevent misinterpretation of the work.
		- Reinforce the execution-architecture focus of gPCD.

		Contains:
		- gPCD does not introduce a new occupancy structure.
		- gPCD does not introduce new asymptotic complexity results.
		- Existing collision-detection techniques are retained.
		- The contribution lies in architectural organization.
		- Existing techniques are integrated into a unified GPU-driven framework.
		- The novelty is execution architecture rather than methodology.

		Presentation Bullets:
		- Not a new data structure.
		- Not a new complexity improvement.
		- Existing techniques are retained.
		- Novelty lies in execution architecture.
		- Unified GPU-driven framework.



   
3. Methodology


	Figure: Overall Framework Architecture

		Purpose:
		- Provide a high-level view of gPCD.
		- Show the relationship between graphics and compute pipelines.
		- Orient the reader before implementation details.

		Supports:
		- gPCD is an execution architecture.
		- Occupancy generation occurs in the graphics pipeline.
		- Narrow-phase evaluation occurs in the compute pipeline.
		- Rendering and occupancy construction proceed concurrently.

		Presentation Bullets:
		- Graphics-pipeline occupancy generation.
		- PCS generation.
		- Compute-based narrow phase.
		- Concurrent rendering and collision detection.
   
 

	Figure: Cell Spans

		Purpose:
		- Illustrate how particle geometry spans cells.
		- Show the relationship between particle size and cell occupancy.
		- Motivate occupancy-list sizing and bounded occupancy generation.
		- Provide geometric justification for the broad-phase design.

		Supports:
		- Particles are represented by bounding volumes.
		- Particle corners determine occupied cells.
		- A particle may intersect multiple neighboring cells.
		- Cell occupancy can be determined geometrically.
		- Occupancy generation is suitable for parallel execution.
		- Occupancy-list capacity can be bounded from particle geometry.

		Presentation Bullets:
		- Particle-to-cell relationships.
		- Corner-based cell assignment.
		- Multi-cell occupancy.
		- Geometric basis for occupancy generation.  
   

	Paragraph 1

		Purpose:
		- Introduce the gPCD framework.
		- Connect the Methodology section to the architectural gap identified in the Background section.
		- State the central architectural idea.
		- Provide a high-level description of the framework before implementation details.

		Contains:
		- gPCD addresses the architectural gap identified in the Background section.
		- Broad-phase occupancy generation is integrated directly into the graphics pipeline.
		- Occupancy construction and rendering proceed concurrently.
		- The framework produces the potentially colliding set (PCS).
		- The PCS is used during narrow-phase evaluation.
		- The framework unifies rendering and collision-detection activities within a GPU-driven architecture.

		Presentation Bullets:
		- gPCD addresses the architectural gap.
		- Graphics-pipeline occupancy generation.
		- Concurrent occupancy construction and rendering.
		- PCS generation for narrow-phase evaluation.
   
3.1 Framework Overview

	Paragraph 1

		Purpose:
		- Explain the architecture figure.
		- Contrast conventional and gPCD organizations.
		- Highlight the architectural difference between existing methods and gPCD.
		- Introduce concurrent occupancy construction and rendering.

		Contains:
		- Figure (a) shows conventional GPU-based collision-detection frameworks.
		- Conventional frameworks separate rendering and collision detection.
		- Figure (b) shows the gPCD architecture.
		- gPCD integrates occupancy generation into the graphics pipeline.
		- Occupancy construction and rendering proceed concurrently.
		- The framework generates the PCS for narrow-phase evaluation.

		Presentation Bullets:
		- Conventional: separate rendering and collision detection.
		- gPCD: graphics-pipeline occupancy generation.
		- Concurrent rendering and occupancy construction.
		- PCS generation integrated into the architecture.  
   
   
3.2 Duplicate Handling

	Paragraph 1

		Purpose:
		- Define the two forms of duplication present in gPCD.
		- Distinguish occupancy duplication from collision duplication.
		- Introduce the duplicate-handling problem.
		- Connect the discussion to the Cell Spans figure.

		Contains:
		- Two forms of duplication can occur in gPCD.
		- Corner duplication occurs during occupancy generation.
		- Multiple AABB corners from the same particle may map to the same cell.
		- Collision duplication occurs during PCS traversal.
		- The same source-target interaction may be encountered through multiple cells.
		- Figure~\ref{fig:CellSpans} illustrates both duplication mechanisms.

		Presentation Bullets:
		- Two duplication mechanisms exist.
		- Corner duplication:
			- Multiple corners → same cell.
		- Collision duplication:
			- Same pair encountered multiple times.
		- Duplicate handling is required.  


	Paragraph 2

		Purpose:
		- Explain how duplicate occupancy entries are prevented.
		- Explain how duplicate collision evaluations are prevented.
		- Associate each duplication type with its corresponding solution.
		- Establish correctness of PCS generation and collision evaluation.

		Contains:
		- Corner duplicates are removed during occupancy construction.
		- Duplicate particle insertions into a cell occupancy list are prevented.
		- Collision duplicates are removed during narrow-phase evaluation.
		- Source-owned collision tracking is used.
		- Each source-target interaction is evaluated at most once.
		- Duplicate suppression preserves correctness and efficiency.

		Presentation Bullets:
		- Corner duplicates removed during occupancy construction.
		- Prevents duplicate cell occupancy entries.
		- Collision duplicates removed during narrow phase.
		- Source-owned collision tracking.
		- Each collision pair evaluated once.


3.3 Graphics-Pipeline Occupancy Generation

	Paragraph 1

		Purpose:
		- Introduce graphics-pipeline occupancy generation.
		- Explain where broad-phase construction occurs.
		- Describe the role of individual particles during occupancy generation.
		- Connect occupancy generation to PCS construction.

		Contains:
		- Broad-phase occupancy generation is performed in the graphics pipeline.
		- Occupancy generation occurs in the vertex stage.
		- Each particle independently determines its occupancy contribution.
		- Each particle determines the cell locations required for broad-phase construction.
		- Occupancy information is stored in the cell-array structure.
		- Cell occupancy information is later used to construct candidate collision pairs.
		- Candidate pairs are passed to the narrow phase.

		Presentation Bullets:
		- Occupancy generation occurs in the graphics pipeline.
		- Implemented in the vertex stage.
		- Particles independently determine occupied cells.
		- Occupancy stored in cell array.
		- Cell array used to construct PCS candidates.

3.4 Corner-Based Occupancy Construction

	Paragraph 1

		Purpose:
		- Introduce the geometric representation used for occupancy generation.
		- Explain how occupied cells are identified.
		- Establish the role of AABB corners in broad-phase construction.
		- Connect particle geometry to cell-array occupancy.

		Contains:
		- Particles are represented by axis-aligned bounding boxes (AABBs).
		- Occupancy generation evaluates the eight AABB corners.
		- Corner locations are mapped to cell locations.
		- Occupied cells are determined from corner-cell intersections.
		- Corner locations define the occupancy contribution of a particle.
		- Cell occupancy is derived directly from particle geometry.

		Presentation Bullets:
		- Particles represented by AABBs.
		- Eight corners evaluated per particle.
		- Corners mapped to cell locations.
		- Corner locations define cell occupancy.



	Paragraph 2

		Purpose:
		- Define the corner-generation procedure mathematically.
		- Show how AABB corners are constructed from particle position and radius.
		- Connect particle geometry to cell occupancy generation.
		- Establish the inputs required for cell-coordinate determination.

		Contains:
		- Particle position is represented by (x,y,z).
		- Particle radius is represented by r.
		- Eight AABB corners are generated using -+r offsets.
		- Each corner corresponds to one of eight corner identifiers.
		- Corner positions are computed directly from particle geometry.
		- Corner locations are converted into cell coordinates.
		- Cell coordinates determine occupancy contributions.

		Presentation Bullets:
		- Corner location:
			x_c = x + (+-r, +-r, +-r)
		- Eight corners per particle.
		- Corner identifiers c in {0,...,7}.
		- Corners converted to cell coordinates.
		- Cell coordinates define occupancy.


	Paragraph 3

		Purpose:
		- Explain the computational consequence of corner-based occupancy generation.
		- Establish bounded work per particle.
		- Connect particle geometry to computational cost.
		- Motivate the scalability of occupancy generation.

		Contains:
		- Occupancy generation uses a fixed set of eight corners.
		- Each particle performs a bounded number of occupancy evaluations.
		- Occupancy work does not depend on total particle count.
		- Per-particle occupancy generation cost is constant.
		- Occupancy generation requires O(1) work per particle.
		- The total occupancy-generation workload scales linearly with particle count.

		Presentation Bullets:
		- Fixed number of corner evaluations.
		- Bounded work per particle.
		- O(1) occupancy generation per particle.
		- Linear scaling with particle count.



	Paragraph 4

		Purpose:
		- State the particle-size constraint used by the framework.
		- Justify why eight-corner occupancy determination is sufficient.
		- Establish the maximum number of cells a particle may occupy.
		- Show that occupancy-generation cost is independent of particle size within the admissible range.

		Contains:
		- Particle diameter must not exceed cell width.
		- Under this constraint, a particle intersects at most eight neighboring cells.
		- Occupancy can be determined entirely from the eight AABB corners.
		- No additional occupancy evaluations are required for admissible particle sizes.
		- Particle size may vary within the admissible range.
		- The occupancy-generation procedure remains unchanged.
		- Per-particle occupancy cost remains bounded.

		Presentation Bullets:
		- Constraint: particle diameter <= cell width.
		- Maximum occupancy span: 8 cells.
		- Eight corners fully determine occupancy.
		- Particle size variation does not increase occupancy work.
		- Constant-cost occupancy generation preserved.


3.5 Cell Addressing and Flattening

	Paragraph 1

		Purpose:
		- Define the mapping from corner positions to cell coordinates.
		- Establish the coordinate system used by the cell array.
		- Explain how floating-point geometry is converted into discrete cell locations.
		- Provide the first step of the cell-addressing process.

		Contains:
		- Each corner location is mapped to a cell coordinate.
		- Corner locations are represented in floating-point coordinates.
		- Cell locations are represented by unsigned integer coordinates.
		- Cell coordinates are obtained by rounding corner positions.
		- Rounding maps each corner to the nearest cell center.
		- Cell coordinates provide the basis for occupancy insertion.

		Presentation Bullets:
		- Corner locations use floating-point coordinates.
		- Cell locations use integer coordinates.
		- Cell coordinates obtained by rounding.
		- Corner-to-cell mapping enables occupancy construction.



	Paragraph 2

		Purpose:
		- Define the valid occupancy region of the cell array.
		- Establish the indexing convention used by the framework.
		- Explain the reservation of zero-valued coordinate locations.
		- Define the admissible cell-coordinate domain.

		Contains:
		- Cell coordinates use zero-based indexing.
		- Cells containing zero-valued coordinate components are reserved.
		- Reserved cells are excluded from occupancy generation.
		- Valid occupancy cells satisfy:
			X > 0
			Y > 0
			Z > 0
		- Occupancy generation operates only within the admissible cell region.
		- Cell-coordinate validity is enforced before occupancy insertion.

		Presentation Bullets:
		- Zero-based indexing.
		- Zero-coordinate cells are reserved.
		- Reserved cells are excluded from occupancy generation.
		- Valid cells satisfy:
			X > 0, Y > 0, Z > 0.



	Paragraph 3

		Purpose:
		- Justify the reserved-cell convention.
		- Explain how boundary conditions are handled during occupancy generation.
		- Prevent invalid cell addressing.
		- Improve robustness of corner-based occupancy construction.

		Contains:
		- Reserved cells create a guard layer around the active domain.
		- The guard layer surrounds valid occupancy cells.
		- Particles near domain boundaries may generate corner locations outside the active region.
		- The guard layer prevents negative cell coordinates.
		- Boundary particles can be processed without special-case logic.
		- Occupancy generation remains robust near domain boundaries.

		Presentation Bullets:
		- Reserved cells form a guard layer.
		- Protects occupancy generation near boundaries.
		- Prevents negative cell coordinates.
		- Eliminates boundary-addressing issues.



	Paragraph 4

		Purpose:
		- Define the mapping from cell coordinates to storage locations.
		- Convert three-dimensional cell coordinates into a one-dimensional index.
		- Enable efficient occupancy-list storage and retrieval.
		- Establish the addressing mechanism used by the cell array.

		Contains:
		- Cell coordinates are converted into a flattened index.
		- The cell array is stored as a one-dimensional structure.
		- W, H, and D represent cell-array dimensions.
		- The flattening operation maps (X,Y,Z) to a unique index I.
		- The flattened index identifies the occupancy-list location.
		- Occupancy insertion and retrieval use the flattened index.

		Presentation Bullets:
		- Convert (X,Y,Z) to a unique index.
		- Flattening enables 1D storage.
		- Cell-array dimensions:
			W, H, D.
		- Flattened index:
			I = X + W(Y + HZ).
		- Supports efficient occupancy access.


	Paragraph 5

		Purpose:
		- Explain the role of the flattened index within the framework.
		- Justify the use of a flattened representation.
		- Describe how the index is used throughout the collision-detection pipeline.
		- Highlight storage and addressing benefits.

		Contains:
		- The flattened index is the primary cell identifier.
		- The index directly addresses occupancy lists.
		- Flattening reduces storage requirements.
		- A three-component coordinate tuple is replaced by a single integer.
		- The same identifier is reused throughout the pipeline.
		- Occupancy construction uses the flattened index.
		- Candidate generation uses the flattened index.
		- Narrow-phase evaluation uses the flattened index.
		- Cell coordinates do not need to be reconstructed.

		Presentation Bullets:
		- Flattened index is the primary cell identifier.
		- Direct occupancy-list addressing.
		- Reduced storage requirements.
		- Single identifier reused throughout the pipeline.
		- No coordinate reconstruction required.

	Paragraph 6

		Purpose:
		- Explain how concurrent occupancy insertion is performed.
		- Introduce thread-safe occupancy-list updates.
		- Describe the role of atomic operations.
		- Ensure correctness during massively parallel occupancy generation.

		Contains:
		- Multiple particles may target the same occupancy list concurrently.
		- Shared occupancy lists require coordinated insertion.
		- Atomic reservation operations are used.
		- Atomic operations allocate unique occupancy-list slots.
		- Each insertion receives a distinct storage location.
		- Thread-safe occupancy generation is maintained.
		- Concurrent occupancy construction is supported.

		Presentation Bullets:
		- Occupancy lists are shared resources.
		- Multiple particles may insert concurrently.
		- Atomic reservation obtains a unique slot.
		- Thread-safe occupancy construction.
		- Supports massively parallel execution.


3.6 Cell Array Construction

	Paragraph 1

		Purpose:
		- Explain how duplicate corner references affect occupancy storage.
		- Describe duplicate suppression during cell-array construction.
		- Establish how occupancy-list capacity is determined.
		- Connect duplicate handling to memory requirements.

		Contains:
		- Multiple AABB corners from the same particle may map to a common cell.
		- Duplicate corner references are suppressed during occupancy generation.
		- Only one particle identifier is stored per particle-cell pair.
		- Occupancy lists store unique particle references.
		- Occupancy-list capacity depends on unique particle contributors.
		- Capacity is not determined by the total number of corner insertions.
		- Duplicate suppression reduces storage requirements.

		Presentation Bullets:
		- Multiple corners may target the same cell.
		- Duplicate corner references are removed.
		- One particle identifier per cell.
		- Occupancy capacity based on unique particles.
		- Capacity independent of duplicate corner insertions.

	Paragraph 2

		Purpose:
		- Derive the maximum occupancy-list capacity.
		- Establish a sizing rule for heterogeneous particle systems.
		- Identify the governing particle size.
		- Provide a memory-allocation model for the cell array.

		Contains:
		- Occupancy capacity depends on particle size.
		- The smallest particle radius governs maximum occupancy.
		- r_min is defined as the minimum particle radius in the simulation.
		- Smaller particles permit higher local packing densities.
		- The uniform cell side length is represented by h.
		- A closed-form expression is provided for maximum cell occupancy.
		- The expression defines the required occupancy-list capacity.
		- Occupancy-list sizing supports heterogeneous particle systems.

		Presentation Bullets:
		- Maximum occupancy determined by smallest particle.
		- Define:
			r_min = min(r_i).
		- Smaller particles create worst-case occupancy.
		- Maximum occupancy:
			N_cell^3D.
		- Provides occupancy-list sizing rule.



	Paragraph 3

		Purpose:
		- Explain how the occupancy-capacity expression is used in practice.
		- Connect the occupancy model to memory allocation.
		- Justify sizing occupancy lists using the smallest particle size.
		- Explain why larger particles cannot increase maximum occupancy.

		Contains:
		- Occupancy-list capacity can be determined during initialization.
		- Capacity is based on the smallest allowable particle diameter.
		- Memory allocation can be performed before simulation begins.
		- Larger particles occupy more volume.
		- Larger particles reduce local packing density.
		- Larger particles cannot increase the maximum number of unique particle identifiers within a cell.
		- The smallest particle defines the worst-case occupancy condition.

		Presentation Bullets:
		- Occupancy capacity determined during initialization.
		- Smallest particle diameter defines worst-case occupancy.
		- Memory allocation performed a priori.
		- Larger particles cannot increase occupancy capacity.
		- Single sizing rule for heterogeneous systems.


3.7 Potentially Colliding Set Construction

	Paragraph 1

		Purpose:
		- Define the potentially colliding set (PCS).
		- Explain the relationship between the cell array and the PCS.
		- Show how occupancy generation produces collision candidates.
		- Eliminate the need for a separate candidate-construction stage.

		Contains:
		- The populated cell array forms the PCS.
		- Occupied cells contain particle identifiers.
		- Particle identifiers are associated with mapped corner locations.
		- Particles sharing occupancy cells become collision candidates.
		- Occupancy generation directly constructs the PCS.
		- No separate candidate-construction stage is required.
		- Broad-phase construction and PCS construction occur simultaneously.

		Presentation Bullets:
		- Populated cell array = PCS.
		- Occupied cells contain particle identifiers.
		- Shared occupancy implies candidate interaction.
		- PCS generated directly during occupancy construction.
		- No separate candidate-generation stage required.


3.8 Narrow-Phase Evaluation

	Paragraph 1

		Purpose:
		- Introduce the data structure used during narrow-phase evaluation.
		- Explain how occupancy information is carried forward from the broad phase.
		- Connect occupancy construction to collision evaluation.
		- Establish the per-particle information required for PCS traversal.

		Contains:
		- Each particle stores an array of eight corner-cell identifiers.
		- Each entry corresponds to one AABB corner.
		- Entries contain flattened cell indices.
		- Cell identifiers are generated during occupancy construction.
		- Occupancy information is retained after broad-phase generation.
		- Narrow-phase evaluation uses the stored corner-cell identifiers.
		- Broad-phase and narrow-phase stages are linked through the cornerCell array.

		Presentation Bullets:
		- Each particle stores cornerCell[8].
		- One entry per AABB corner.
		- Entries contain flattened cell indices.
		- Generated during occupancy construction.
		- Reused during narrow-phase evaluation.




	Paragraph 2

		Purpose:
		- Explain how collision candidates are retrieved.
		- Describe source-owned PCS traversal.
		- Eliminate the need for centralized collision-pair construction.
		- Connect occupancy lists to contact evaluation.

		Contains:
		- Each source particle independently performs candidate retrieval.
		- The source particle traverses its cornerCell array.
		- Each corner-cell identifier references an occupancy list.
		- Candidate neighbors are obtained directly from occupancy lists.
		- No centralized collision-pair list is constructed.
		- Collision candidates are generated on demand.
		- The source particle evaluates contact against referenced particles.
		- Candidate retrieval and collision evaluation are source owned.

		Presentation Bullets:
		- Source particle traverses cornerCell[8].
		- Occupancy lists provide candidate neighbors.
		- No centralized collision-pair list.
		- Candidates generated on demand.
		- Source-owned collision evaluation.



	Paragraph 3

		Purpose:
		- Explain duplicate-collision suppression during narrow-phase evaluation.
		- Describe the source-owned collision-tracking mechanism.
		- Ensure each source-target interaction is evaluated at most once.
		- Preserve correctness during PCS traversal.

		Contains:
		- Each particle maintains a collision list.
		- The collision list stores previously evaluated target identifiers.
		- Collision tracking is source owned.
		- When a collision is detected, the target identifier is recorded.
		- Future evaluations consult the collision list.
		- Previously evaluated source-target pairs are discarded.
		- Duplicate collision evaluations are prevented.
		- Each collision pair is evaluated at most once per evaluation step.

		Presentation Bullets:
		- Source-owned collision list.
		- Stores previously evaluated targets.
		- Collision registration updates the list.
		- Duplicate source-target evaluations discarded.
		- One evaluation per collision pair.



	Paragraph 4

		Purpose:
		- Explain the benefits of source-owned candidate evaluation.
		- Justify the narrow-phase execution architecture.
		- Reduce synchronization requirements.
		- Demonstrate suitability for massively parallel execution.

		Contains:
		- Candidate discovery is performed independently by each source particle.
		- Candidate ownership remains local to the source particle.
		- Collision evaluation is source owned.
		- No globally shared collision-pair list is required.
		- Synchronization requirements are reduced.
		- Parallel execution is simplified.
		- The organization is well suited to massively parallel hardware.
		- Collision management remains decentralized.

		Presentation Bullets:
		- Source-owned candidate discovery.
		- Local candidate ownership.
		- No global collision-pair structure.
		- Reduced synchronization.
		- Massively parallel execution.

4. Complexity Analysis

		Table: Complexity Comparison

		Purpose:
		- Compare the theoretical complexity of gPCD with representative
		  broad-phase collision-detection methodologies.
		- Position gPCD within the existing collision-detection literature.
		- Provide context for the theoretical scaling analysis.
		- Highlight differences between methodologies and execution architectures.

		Supports:
		- gPCD achieves competitive theoretical complexity.
		- Architectural innovation does not require asymptotic-complexity improvements.
		- Broad-phase methodologies exhibit differing computational characteristics.
		- The contribution of gPCD is architectural rather than asymptotic.

		Presentation Bullets:
		- Complexity comparison with representative methods.
		- Broad-phase computational complexity.
		- Relative scalability.
		- Architectural innovation without asymptotic penalty.

	Paragraph 1

		Purpose:
		- Introduce the complexity-analysis section.
		- Identify the dominant computational components of the framework.
		- Focus the analysis on occupancy generation and collision evaluation.
		- Establish the scope of the theoretical analysis.

		Contains:
		- The computational complexity of gPCD is analyzed.
		- Broad-phase occupancy generation is a primary contributor to cost.
		- Narrow-phase collision evaluation is a primary contributor to cost.
		- These stages determine overall framework complexity.
		- Complexity analysis focuses on the dominant computational operations.

		Presentation Bullets:
		- Complexity driven by:
			* Broad-phase occupancy generation
			* Narrow-phase collision evaluation
		- Dominant framework costs analyzed.

4.1 Occupancy Generation

	Paragraph 1

		Purpose:
		- Derive the computational complexity of occupancy generation.
		- Formalize the bounded-work argument from the methodology.
		- Establish the broad-phase scaling behavior.
		- Provide the first theoretical complexity result for gPCD.

		Contains:
		- Each particle performs a fixed number of occupancy evaluations.
		- Occupancy work is independent of total particle count.
		- c_occ represents the bounded occupancy cost per particle.
		- Occupancy-generation cost grows linearly with particle count.
		- Broad-phase occupancy generation has O(N) complexity.
		- The broad phase exhibits constant work per particle.

		Presentation Bullets:
		- Fixed occupancy evaluations per particle.
		- Bounded occupancy cost c_occ.
		- Linear growth with particle count.
		- Broad-phase occupancy generation:O(N).

	Paragraph 2

		Purpose:
		- Establish the cost of cell-array access operations.
		- Show that addressing does not increase occupancy-generation complexity.
		- Connect flattened indexing to computational efficiency.
		- Reinforce the bounded-work assumption.

		Contains:
		- Occupancy lists are accessed through flattened indices.
		- Cell addressing is performed in constant time.
		- Direct indexing eliminates search overhead.
		- Occupancy insertion requires bounded work.
		- Occupancy retrieval requires bounded work.
		- Fixed occupancy-width assumptions are maintained.
		- Addressing and storage operations preserve O(1) per-particle behavior.
		- Occupancy-generation complexity remains O(N).

		Presentation Bullets:
		- Direct access via flattened indices.
		- Constant-time cell lookup.
		- Bounded insertion cost.
		- Bounded retrieval cost.
		- Addressing preserves O(N) occupancy generation.

4.2 Narrow-Phase Evaluation

	Paragraph 1

	Purpose:
	- Derive the computational complexity of narrow-phase evaluation.
	- Define the quantity governing narrow-phase cost.
	- Establish the scaling behavior of collision evaluation.
	- Separate occupancy-generation complexity from interaction-evaluation complexity.

	Contains:
	- Source particles traverse occupancy lists associated with corner cells.
	- Narrow-phase cost depends on the number of interaction evaluations.
	- K represents the total number of candidate interactions evaluated.
	- Each interaction requires bounded work.
	- Narrow-phase cost grows linearly with K.
	- Narrow-phase evaluation has O(K) complexity.
	- Collision-processing cost is proportional to the number of evaluated interactions.

	Presentation Bullets:
	- Source particles traverse occupancy lists.
	- K = total interaction evaluations.
	- Constant work per interaction.
	- Narrow-phase complexity:O(K).



	Paragraph 2

		Purpose:
		- Interpret the narrow-phase complexity result.
		- Identify the quantity that governs collision-evaluation cost.
		- Distinguish interaction count from particle count.
		- Connect PCS construction to narrow-phase workload.

		Contains:
		- Narrow-phase complexity depends on candidate interactions.
		- Candidate interactions are represented by the PCS.
		- Interaction count is the primary driver of narrow-phase cost.
		- Particle count alone does not determine narrow-phase workload.
		- Different particle distributions may produce different values of K.
		- PCS density influences collision-processing cost.
		- Narrow-phase scaling is governed by interaction density rather than particle count alone.

		Presentation Bullets:
		- Narrow-phase cost depends on PCS interactions.
		- K is more important than N.
		- Particle count alone does not determine workload.
		- Interaction density drives collision-evaluation cost.


4.3 Storage Requirements

	Paragraph 1

		Purpose:
		- Derive the storage requirements of the framework.
		- Separate particle-storage costs from cell-array storage costs.
		- Establish the memory-scaling behavior of gPCD.
		- Provide the storage-complexity result.

		Contains:
		- Particle storage scales linearly with particle count.
		- M_particle = O(N).
		- Cell-array storage depends on array dimensions and occupancy width.
		- M_cell = O(S^3 W).
		- Occupancy width W is fixed.
		- Cell-array storage reduces to O(S^3).
		- Total storage is the sum of particle and cell-array storage.
		- M_gPCD = O(N + S^3 W).
		- For fixed occupancy width:
			M_gPCD = O(N + S^3).

		Presentation Bullets:
		- Particle storage:
			O(N).
		- Cell-array storage:
			O(S^3 W).
		- Fixed occupancy width:
			O(S^3).
		- Total storage:
			O(N + S^3).

5. Comparative Analysis

	Table: Comparative Complexity

		Purpose:
		- Compare gPCD with representative collision-detection methodologies.
		- Position the framework relative to existing approaches.
		- Contrast methodological and architectural contributions.
		- Provide context for interpreting the complexity results.

		Supports:
		- gPCD exhibits competitive theoretical complexity.
		- gPCD does not require a new asymptotic-complexity class.
		- Direct-insertion spatial subdivision approaches share similar broad-phase scaling.
		- The primary distinction of gPCD is architectural.
		- Architectural innovation can occur without asymptotic-complexity improvement.

		Presentation Bullets:
		- Complexity comparison with representative methods.
		- Comparable broad-phase scaling.
		- No new asymptotic complexity class.
		- Architectural differentiation.

	Table: Historical Comparison

		Purpose:
		- Place gPCD within the historical development of particle simulation.
		- Compare representative particle counts and TPP values.
		- Provide context for interpreting reported performance.
		- Illustrate how simulation scale has evolved across studies.

		Supports:
		- Published studies span a wide range of particle counts.
		- Published studies span a wide range of TPP values.
		- Direct performance comparisons are difficult due to differing
		  hardware, algorithms, and benchmark conditions.
		- gPCD supports large particle counts on a single CPU-GPU system.
		- gPCD exhibits competitive time-per-particle performance.

		Presentation Bullets:
		- Historical particle-count comparison.
		- Historical TPP comparison.
		- Literature context.
		- gPCD positioning.

	

	Paragraph 1

		Purpose:
		- Clarify the nature of the gPCD contribution.
		- Distinguish architectural innovation from asymptotic-complexity improvements.
		- Position gPCD among direct-insertion spatial subdivision methods.
		- Summarize the outcome of the comparative analysis.

		Contains:
		- gPCD does not introduce a new asymptotic complexity class.
		- Broad-phase occupancy generation remains linear.
		- Similar broad-phase complexity exists in other direct-insertion methods.
		- The contribution is architectural rather than asymptotic.
		- Graphics-pipeline occupancy generation is a distinguishing feature.
		- Duplicate-aware cell construction is a distinguishing feature.
		- Source-owned collision evaluation is a distinguishing feature.
		- These elements are unified within a GPU-driven framework.

		Presentation Bullets:
		- No new asymptotic complexity class.
		- Broad phase remains O(N).
		- Contribution is architectural.
		- Graphics-pipeline occupancy generation.
		- Duplicate-aware cell construction.
		- Source-owned collision evaluation.
		- Unified GPU-driven framework.


	Paragraph 2

		Purpose:
		- Interpret the historical-comparison table.
		- Caution against direct performance ranking.
		- Position gPCD relative to prior work.
		- Highlight achieved particle count and TPP.

		Contains:
		- Representative particle counts are summarized.
		- Representative TPP values are summarized.
		- Benchmark conditions differ substantially across studies.
		- Hardware platforms differ across studies.
		- Reported operations differ across studies.
		- Results should be interpreted as historical context.
		- The comparison is not a strict performance ranking.
		- gPCD supports more than 11 million particles.
		- gPCD achieves approximately 7.62 ns per particle.
		- Results are obtained on a single CPU-GPU system.

		Presentation Bullets:
		- Literature comparisons require caution.
		- Different hardware and benchmarks.
		- Historical context rather than ranking.
		- gPCD:
			• >11 million particles
			• 7.62 ns/particle
			• Single CPU-GPU system


6. Experimental Setup

	Paragraph 1

		Purpose:
		- Introduce the experimental-evaluation section.
		- Define the scope of the benchmark methodology.
		- Identify the elements required to reproduce the evaluation.
		- Prepare the reader for the performance results.

		Contains:
		- The evaluation methodology is described.
		- Hardware platform is specified.
		- Benchmark configurations are specified.
		- Performance metrics are specified.
		- These components are used to evaluate gPCD.
		- The section establishes the experimental basis for the Results section.

		Presentation Bullets:
		- Experimental methodology.
		- Hardware platform.
		- Benchmark configurations.
		- Performance metrics.
		- Evaluation framework.


6.1 Hardware Platform

	Paragraph 1

		Purpose:
		- Define the hardware platform used for evaluation.
		- Establish the computational resources available to gPCD.
		- Provide context for interpreting performance results.
		- Position the benchmark system relative to practical deployment environments.

		Contains:
		- Evaluation performed on a Lenovo ThinkPad P16 Gen 2.
		- CPU:
			Intel Core i9-13980HX.
		- GPU:
			NVIDIA RTX 3500 Ada Generation.
		- GPU contains 5,120 CUDA cores.
		- GPU provides 12 GB ECC GDDR6 memory.
		- GPU uses a 192-bit memory interface.
		- GPU bandwidth is approximately 432 GB/s.
		- System represents a modern mobile workstation.
		- Hardware serves as the evaluation baseline.

		Presentation Bullets:
		- Lenovo ThinkPad P16 Gen 2.
		- Intel Core i9-13980HX.
		- NVIDIA RTX 3500 Ada.
		- 5,120 CUDA cores.
		- 12 GB ECC GDDR6.
		- 432 GB/s memory bandwidth.
		- Mobile workstation platform.

	Paragraph 2

		Purpose:
		- Describe the software implementation of gPCD.
		- Explain how framework operations are mapped onto GPU execution resources.
		- Connect the experimental implementation to the methodology.
		- Establish the execution environment used during benchmarking.

		Contains:
		- gPCD is implemented using Vulkan.
		- Broad-phase occupancy generation executes in the graphics pipeline.
		- Rendering executes in the graphics pipeline.
		- Narrow-phase evaluation executes in compute shaders.
		- Graphics and compute pipelines are both utilized.
		- The implementation follows the proposed GPU-driven architecture.

		Presentation Bullets:
		- Vulkan implementation.
		- Graphics-pipeline occupancy generation.
		- Graphics-pipeline rendering.
		- Compute-shader narrow phase.
		- GPU-driven execution architecture.
	
6.2 Maximum Particle Capacity

	Paragraph 1

		Purpose:
		- Define the theoretical upper bound on particle count.
		- Relate simulation capacity to graphics-device memory limits.
		- Establish a hardware-dependent particle-count constraint.
		- Provide a sizing model for maximum simulation scale.

		Contains:
		- Particle count is limited by storage-buffer capacity.
		- M_max represents the maximum addressable storage-buffer range.
		- S_p represents the memory footprint of a particle.
		- Maximum particle count depends on available buffer storage.
		- N_max is computed from the ratio M_max / S_p.
		- Larger particle structures reduce maximum capacity.
		- Smaller particle structures increase maximum capacity.
		- The capacity limit is hardware dependent.

		Presentation Bullets:
		- Maximum capacity constrained by storage-buffer limits.
		- Define:
			• M_max = maximum storage-buffer range
			• S_p = particle size
		- Maximum particle count:
			N_max = floor(M_max / S_p)
		- Capacity depends on particle footprint.
		- Capacity depends on hardware.



	Paragraph 2

		Purpose:
		- Explain how the capacity model was applied during testing.
		- Connect theoretical capacity limits to benchmark selection.
		- Describe the effect of particle-structure size on simulation scale.
		- Demonstrate that all experiments respected hardware constraints.

		Contains:
		- The capacity equation was used to select benchmark sizes.
		- Tested particle counts were constrained by storage-buffer limits.
		- Increasing particle size reduces maximum achievable particle count.
		- Particle capacity scales inversely with particle memory footprint.
		- Reducing particle size increases achievable simulation scale.
		- Maximum problem size depends on storage efficiency.
		- All benchmark cases remained below the device storage-buffer limit.
		- Experimental configurations were selected to ensure valid execution.

		Presentation Bullets:
		- Capacity equation used to select benchmarks.
		- Larger particle structures reduce maximum capacity.
		- Smaller particle structures increase maximum capacity.
		- Maximum problem size depends on particle memory footprint.
		- All tests remained within hardware limits.



6.3 Verification and Performance Evaluation

	Paragraph 1

		Purpose:
		- Explain the evaluation methodology used for gPCD.
		- Separate correctness validation from performance measurement.
		- Establish confidence in reported benchmark results.
		- Justify the use of verification mode.

		Contains:
		- gPCD performs exact collision detection.
		- Verification and performance evaluations are separate stages.
		- Each test case is first executed in verification mode.
		- Verification validates particle counts.
		- Verification validates occupancy generation.
		- Verification validates collision counts.
		- Verification ensures all particles are processed correctly.
		- Verification ensures collision events are detected and counted correctly.
		- Correctness is established before performance evaluation.

		Presentation Bullets:
		- Exact collision detection framework.
		- Verification and performance measured separately.
		- Verification validates:
			* Particle counts
			* Occupancy generation
			* Collision counts
		- Correctness established before benchmarking.



	Paragraph 2

		Purpose:
		- Describe the performance-evaluation stage.
		- Explain when performance measurements were collected.
		- Ensure that reported metrics are based on validated results.
		- Establish the credibility of the reported benchmarks.

		Contains:
		- Performance evaluation follows successful verification.
		- The same benchmark configurations are reused.
		- Performance mode measures execution time.
		- Performance mode measures throughput.
		- Performance mode measures time-per-particle scaling.
		- Correctness is established before performance measurement.
		- Reported performance metrics correspond to validated collision-detection results.
		- Benchmark results are based on verified executions.

		Presentation Bullets:
		- Performance mode executed after verification.
		- Same benchmark configurations used.
		- Measured:
			• Execution time
			• Throughput
			• Time per particle
		- Performance results based on validated collision detection.


6.3 Benchmark Configuration

	Paragraph 1

		Purpose:
		- Define the benchmark configurations used for evaluation.
		- Describe how test cases were generated.
		- Establish the particle-count range used in scaling studies.
		- Control collision density while varying problem size.

		Contains:
		- Benchmark configurations use randomly distributed particles.
		- Particle count is systematically increased across test cases.
		- Collision fraction is maintained approximately constant.
		- Scaling behavior is evaluated independently of large density variations.
		- Particle counts range from 32 to 11,000,000 particles.
		- Benchmark configurations span multiple orders of magnitude.

		Presentation Bullets:
		- Random particle configurations.
		- Approximately constant collision fraction.
		- Systematic increase in particle count.
		- Range:
			32 to 11,000,000 particles.
		- Multi-order-of-magnitude scaling study.

6.5 Performance Metrics

	Paragraph 1

		Purpose:
		- Define the performance metrics used to evaluate gPCD.
		- Establish the quantities reported in the Results section.
		- Separate graphics, compute, and total execution costs.
		- Introduce throughput as a primary performance metric.

		Contains:
		- Graphics execution time is measured.
		- Compute execution time is measured.
		- Total execution time is measured.
		- Throughput is measured.
		- Time-per-particle is measured.
		- Metrics capture both execution cost and processing efficiency.

		Presentation Bullets:
		- Metrics:
			• Graphics time (t_g)
			• Compute time (t_c)
			• Total time (t_t)
			• Throughput
			• Time per particle



	Paragraph 2

		Purpose:
		- Define throughput mathematically.
		- Relate particle count to execution time.
		- Establish the primary measure of processing rate.
		- Provide a basis for throughput comparisons.

		Contains:
		- Throughput is particles processed per second.
		- Throughput depends on particle count.
		- Throughput depends on total execution time.
		- Throughput is defined as:
			T = N / t_t.
		- Higher throughput indicates greater processing capacity.
		- Throughput is used to evaluate scalability and efficiency.

		Presentation Bullets:
		- Throughput = particles per second.
		- T = N / t_t.
		- Combines particle count and execution time.
		- Measures processing capacity.




	Paragraph 3

		Purpose:
		- Define time-per-particle (TPP).
		- Introduce a normalized performance metric.
		- Enable comparisons across different particle counts.
		- Provide a measure of computational efficiency.

		Contains:
		- Time-per-particle is computed from total execution time.
		- TPP is defined as:
			t_pp = t_t / N.
		- TPP normalizes execution cost by particle count.
		- TPP allows comparison across benchmark sizes.
		- TPP is independent of absolute particle count.
		- TPP is used to evaluate scaling efficiency.

		Presentation Bullets:
		- Time per particle:
			t_pp = t_t / N.
		- Normalized execution cost.
		- Independent of particle count.
		- Measures scaling efficiency.


7. Experimental Results

	Paragraph 1

		Purpose:
		- Introduce the results section.
		- Define the benchmark conditions used for the reported results.
		- Identify the primary performance measures.
		- Establish the scope of the experimental analysis.

		Contains:
		- Measured gPCD performance is presented.
		- Results use randomly distributed particle configurations.
		- Particle count increases across benchmark cases.
		- Collision fraction remains approximately constant unless noted otherwise.
		- Performance is evaluated using throughput.
		- Performance is evaluated using time-per-particle.
		- Performance is evaluated using scaling analyses.
		- Results characterize framework behavior across the tested particle range.

		Presentation Bullets:
		- Experimental performance results.
		- Random particle configurations.
		- Increasing particle count.
		- Approximately constant collision fraction.
		- Throughput analysis.
		- Time-per-particle analysis.
		- Scaling analysis.

7.1 Memory Locality Benchmark
	Figure: Memory Locality Benchmark

		Purpose:
		- Compare execution times for sequential and random occupancy ordering.
		- Quantify locality-induced performance differences.
		- Measure how the locality effect changes with particle count.

		Supports:
		- Both configurations exhibit approximately linear scaling.
		- Sequential ordering consistently outperforms random ordering.
		- The execution-time difference grows approximately linearly.
		- Difference growth is characterized by slope \PPBQDiffSlope.
		- The observed performance gap arises from memory-locality effects.
		- Algorithmic complexity remains unchanged.

	Paragraph 1

		Purpose:
		- Introduce the memory-locality study.
		- Define the benchmark variable being investigated.
		- Describe the two occupancy-ordering configurations.
		- Establish the basis for the locality comparison.

		Contains:
		- The benchmark evaluates memory-locality effects.
		- Memory locality is treated as an independent variable.
		- Two occupancy configurations are compared.
		- Sequential ordering preserves locality within memory.
		- Random ordering disrupts locality within memory.
		- Both configurations use the same occupancy-generation framework.
		- The benchmark isolates the impact of occupancy ordering.

		Presentation Bullets:
		- Memory-locality sensitivity study.
		- Two occupancy configurations:
			* Sequential ordering
			* Random ordering
		- Locality treated as an independent variable.
		- Isolates memory-ordering effects.

	Paragraph 2

		Purpose:
		- Present the locality-benchmark results.
		- Compare sequential and random occupancy ordering.
		- Introduce the execution-time difference metric.
		- Establish the observed performance trend.

		Contains:
		- Execution time is plotted versus particle count.
		- Sequential ordering produces lower execution times.
		- The performance advantage exists throughout the tested range.
		- Random ordering produces higher execution times.
		- A difference curve is included in the figure.
		- The difference is defined as:
			Random Time - Sequential Time.
		- The figure enables direct visualization of locality effects.
		- Performance differences are quantified across the particle range.

		Presentation Bullets:
		- Execution time versus particle count.
		- Sequential ordering consistently faster.
		- Random ordering consistently slower.
		- Difference curve included.
		- Difference:
			Random - Sequential.



	Paragraph 3

		Purpose:
		- Interpret the locality-benchmark results.
		- Characterize the scaling behavior of both configurations.
		- Quantify the growth of the execution-time difference.
		- Attribute the observed performance gap.

		Contains:
		- Sequential and random configurations exhibit approximately linear scaling.
		- Neither configuration exhibits dominant super-linear behavior.
		- The execution-time difference increases approximately linearly with particle count.
		- A linear regression was performed on the difference curve.
		- The regression slope quantifies the rate at which the performance gap grows.
		- The slope equals \PPBQDiffSlope ms per million particles.
		- The locality penalty grows proportionally with problem size.
		- Both configurations execute the same collision-detection algorithm.
		- Occupancy ordering is the only significant difference.
		- The performance gap is attributed to memory-locality effects.

		Presentation Bullets:
		- Both configurations scale approximately linearly.
		- Difference grows approximately linearly.
		- Difference slope:
			\PPBQDiffSlope ms per million particles.
		- Locality effect increases with scale.
		- Performance gap attributed to memory locality.



	Paragraph 4

		Purpose:
		- Summarize the implications of the locality benchmark.
		- Distinguish algorithmic complexity from practical performance.
		- Explain the role of memory locality in observed execution times.
		- Connect the benchmark results to real-world performance.

		Contains:
		- Memory locality remains an important performance factor.
		- The effect persists despite GPU-resident execution.
		- Memory locality influences execution efficiency.
		- Memory locality does not alter algorithmic complexity.
		- The asymptotic scaling behavior remains unchanged.
		- Memory locality affects implementation-dependent constants.
		- Practical execution time depends on these constant factors.
		- Observed performance differences arise from memory-access behavior rather than computational complexity.

		Presentation Bullets:
		- Memory locality matters.
		- GPU execution does not eliminate locality effects.
		- Complexity unchanged.
		- Constant factors change.
		- Practical performance depends on memory access patterns.

7.2 Performance Overview

	Table: Benchmark Measurements

		Purpose:
		- Present the benchmark data used throughout the Results section.
		- Provide the measured execution statistics for each benchmark case.
		- Serve as the source dataset for all subsequent analyses.
		- Support reproducibility and independent interpretation.

		Supports:
		- Throughput calculations.
		- Time-per-particle calculations.
		- Scaling analyses.
		- Regression analyses.
		- Curvature analyses.
		- Performance-overview conclusions.

		Presentation Bullets:
		- Particle count.
		- Graphics execution time.
		- Compute execution time.
		- Total execution time.
		- Throughput.
		- Time per particle.

	Paragraph 1

		Purpose:
		- Introduce the overall performance results.
		- Reference the summary figure/table containing the principal metrics.
		- Define the benchmark conditions associated with the reported results.
		- Provide context for the detailed analyses that follow.

		Contains:
		- Principal gPCD performance metrics are summarized.
		- Performance results are aggregated into a summary presentation.
		- Benchmark cases use randomly distributed particles.
		- Particle count increases across benchmark configurations.
		- Collision fraction remains approximately constant.
		- Results characterize performance across the tested particle range.
		- The summary provides a high-level view of framework behavior.

		Presentation Bullets:
		- Performance summary.
		- Random particle configurations.
		- Increasing particle count.
		- Approximately constant collision fraction.
		- High-level performance overview.


	Paragraph 2

		Purpose:
		- Present the principal performance achievements of gPCD.
		- Report the largest benchmark successfully processed.
		- Report peak throughput performance.
		- Reference representative execution statistics.

		Contains:
		- gPCD successfully processed the largest benchmark configuration.
		- Maximum tested particle count:
			\APQBRTTABLEALLbka{} particles.
		- Peak throughput:
			\SUMMARYTHROUGHPUT{} million particles/s.
		- Large-scale particle processing was demonstrated.
		- Representative execution statistics are reported.
		- The largest benchmark configuration is summarized in the performance table.
		- The table provides detailed execution metrics for the largest case.

		Presentation Bullets:
		- Largest benchmark:
			\APQBRTTABLEALLbka{} particles.
		- Peak throughput:
			\SUMMARYTHROUGHPUT{} million particles/s.
		- Large-scale processing demonstrated.
		- Largest-case execution statistics reported.	
		


	Paragraph 3

		Purpose:
		- Summarize the overall performance behavior of gPCD.
		- Relate observed behavior to expectations for GPU-resident execution.
		- Establish that performance remained stable across the tested range.
		- Transition to the detailed performance analyses.

		Contains:
		- Performance remained stable across benchmark sizes.
		- Framework behavior was consistent throughout the evaluated range.
		- Throughput-oriented GPU scaling behavior was observed.
		- Results are consistent with expectations for a GPU-resident framework.
		- No unexpected performance degradation was observed.
		- Detailed throughput analysis follows.
		- Detailed time-per-particle analysis follows.
		- Detailed scaling analysis follows.

		Presentation Bullets:
		- Stable performance across benchmark range.
		- Expected GPU-resident behavior.
		- Throughput-oriented scaling.
		- Detailed analysis follows:
			* Throughput
			* Time per particle
			* Scaling behavior		

	Paragraph 4

		Purpose:
		- Identify the benchmark dataset used in the Results section.
		- Reference the source of the reported performance measurements.
		- Establish traceability between analyses and measured data.
		- Provide a common data source for subsequent performance discussions.

		Contains:
		- Benchmark measurements are tabulated.
		- The table contains the data used throughout subsequent analyses.
		- Throughput analysis uses the table data.
		- Time-per-particle analysis uses the table data.
		- Scaling analysis uses the table data.
		- The table serves as the primary benchmark reference.

		Presentation Bullets:
		- Benchmark measurement table.
		- Source data for all performance analyses.
		- Throughput results.
		- Time-per-particle results.
		- Scaling analyses.		
		


	Paragraph 5

		Purpose:
		- Provide a roadmap for the remaining Results section.
		- Identify the detailed analyses that follow.
		- Transition from summary results to supporting evidence.
		- Establish the structure of the subsequent subsections.

		Contains:
		- Throughput behavior will be analyzed.
		- Time-per-particle trends will be analyzed.
		- Scaling characteristics will be analyzed.
		- Collision-fraction sensitivity will be analyzed.
		- Detailed examination follows the performance overview.
		- Subsequent analyses provide supporting evidence for the summary results.

		Presentation Bullets:
		- Detailed throughput analysis.
		- Detailed time-per-particle analysis.
		- Detailed scaling analysis.
		- Detailed collision-fraction study.		
		
		
	7.3 Throughput and Time-per-Particle Behavior

	Paragraph 1

		Purpose:
		- Introduce the primary performance metrics.
		- Explain the relationship between throughput and time-per-particle.
		- Justify the use of both metrics.
		- Establish the interpretation framework for the results.

		Contains:
		- Throughput and time-per-particle are complementary metrics.
		- Throughput measures aggregate processing rate.
		- Throughput characterizes total framework productivity.
		- Time-per-particle normalizes execution cost.
		- Time-per-particle removes the direct influence of particle count.
		- Time-per-particle measures scaling efficiency.
		- Both metrics are required for a complete performance assessment.

		Presentation Bullets:
		- Two complementary metrics:
			• Throughput
			• Time per particle
		- Throughput = processing rate.
		- TPP = scaling efficiency.
		- Both required to characterize performance.		
		
		
7.3.1 Utilization-Limited Regime

	Paragraph 1

		Purpose:
		- Define the low-particle-count operating regime.
		- Explain the initial throughput growth behavior.
		- Relate throughput improvement to GPU utilization.
		- Identify the approximate extent of the utilization-limited region.

		Contains:
		- Low particle counts underutilize available GPU resources.
		- Throughput increases rapidly as particle count increases.
		- Additional particles improve hardware utilization.
		- GPU resources become progressively more occupied.
		- Throughput growth is driven primarily by increasing utilization.
		- The utilization-limited regime extends to approximately 100,000 particles.
		- This regime precedes throughput saturation behavior.

		Presentation Bullets:
		- Low particle-count regime.
		- GPU initially underutilized.
		- Throughput increases rapidly.
		- Increasing utilization drives performance gains.
		- Extends to ~100,000 particles.		
		
7.3.2 Throughput-Transition Regime

	Paragraph 1

		Purpose:
		- Define the intermediate operating regime.
		- Identify the onset of workload-dependent behavior.
		- Explain why throughput growth begins to slow.
		- Establish the particle-count range of the transition region.

		Contains:
		- Throughput behavior changes beyond approximately 100,000 particles.
		- GPU utilization is no longer the sole performance driver.
		- Occupancy-generation workload becomes increasingly significant.
		- Collision-processing workload becomes increasingly significant.
		- Throughput continues to increase.
		- Throughput growth occurs at a reduced rate.
		- The framework transitions from utilization-limited to workload-influenced behavior.
		- The transition regime spans approximately 100,000 to 2,000,000 particles.

		Presentation Bullets:
		- Intermediate performance regime.
		- Begins near 100,000 particles.
		- Occupancy workload increases.
		- Collision workload increases.
		- Throughput growth slows.
		- Extends to ~2,000,000 particles.
		
	7.3.3 Stable Large-Scale Regime

	Paragraph 1

		Purpose:
		- Define the large-scale operating regime.
		- Identify the onset of asymptotic framework behavior.
		- Explain why this region is used for detailed scaling analysis.
		- Establish the particle-count threshold for stable operation.

		Contains:
		- The framework enters a stable operating regime near 2,000,000 particles.
		- Time-per-particle becomes nearly constant.
		- GPU utilization has effectively stabilized.
		- Performance behavior becomes predictable.
		- The regime represents mature large-scale execution.
		- This region is used for subsequent TPP analysis.
		- This region is used for subsequent scaling analysis.
		- The regime approximates the asymptotic behavior of the framework.

		Presentation Bullets:
		- Begins near 2,000,000 particles.
		- Stable large-scale operation.
		- Nearly constant time per particle.
		- Performance behavior stabilizes.
		- Represents asymptotic framework behavior.
		- Used for subsequent scaling analysis.	
				
				
7.3.3 Stable Large-Scale Regime

	Paragraph 1

		Purpose:
		- Define the large-scale operating regime.
		- Explain the relationship between decreasing throughput and stable TPP.
		- Identify the asymptotic operating region.
		- Justify using particle counts above 2,000,000 for TPP analysis.

		Contains:
		- Stable large-scale regime begins near 2,000,000 particles.
		- Throughput continues to decrease gradually as particle count increases.
		- Time-per-particle curves become nearly constant.
		- Nearly constant TPP indicates stable average processing cost per particle.
		- The framework has entered an asymptotic operating region.
		- Particle counts above 2,000,000 are used for the TPP analysis in Figure~\ref{fig:A_PQBR_THROUGHPUT_AND_TPP}(c).

		Presentation Bullets:
		- Stable regime begins near 2 million particles.
		- Throughput gradually decreases with increasing particle count.
		- Time per particle becomes nearly constant.
		- Constant TPP indicates stable average cost.
		- Counts above 2 million used for TPP analysis.		
		
7.4 Scaling Analysis

	Paragraph 1

		Purpose:
		- Introduce the scaling-analysis methodology.
		- Quantify measured framework scaling behavior.
		- Define the regression model used for evaluation.
		- Establish the meaning of the scaling exponent.

		Contains:
		- Scaling behavior is evaluated using log-log regression.
		- Graphics execution times are analyzed.
		- Compute execution times are analyzed.
		- Total execution times are analyzed.
		- A power-law model is fitted to the measured data.
		- The model has the form:
			t = aN^b.
		- t represents execution time.
		- N represents particle count.
		- a is a proportionality constant.
		- b is the scaling exponent.
		- The scaling exponent is the primary quantity used to characterize scaling behavior.

		Presentation Bullets:
		- Log-log regression analysis.
		- Graphics, compute, and total times analyzed.
		- Power-law model:
			t = aN^b
		- a = proportionality constant.
		- b = scaling exponent.
		- Scaling quantified through exponent b.		



	Paragraph 2

		Purpose:
		- Present the measured scaling exponents.
		- Report regression results for graphics, compute, and total execution times.
		- Interpret the meaning of the measured exponents.
		- Quantify observed scaling behavior.

		Contains:
		- Log-log regression results are presented.
		- Graphics execution time exponent is reported.
		- Compute execution time exponent is reported.
		- Total execution time exponent is reported.
		- All exponents remain close to unity.
		- Execution time increases approximately linearly with particle count.
		- Similar scaling behavior is observed across graphics and compute stages.
		- Measured behavior is consistent with near-linear scaling.

		Presentation Bullets:
		- Graphics exponent:
			b_g = \APQBRLOGLOGREGRESSgmsasat{}
		- Compute exponent:
			b_c = \APQBRLOGLOGREGRESScmsasat{}
		- Total exponent:
			b_t = \APQBRLOGLOGREGRESStotasat{}
		- All exponents near 1.
		- Approximately linear scaling observed.		
		
		


	Paragraph 3

		Purpose:
		- Summarize the scaling-analysis results.
		- Relate measured scaling behavior to theoretical predictions.
		- Connect experimental observations to the complexity analysis.
		- Reinforce the near-linear scaling conclusion.

		Contains:
		- Graphics scaling exponent remains close to unity.
		- Compute scaling exponent remains close to unity.
		- Total execution-time exponent remains close to unity.
		- Measured behavior is consistent across framework components.
		- Experimental results support near-linear scaling.
		- Measured behavior agrees with complexity-analysis predictions.
		- Theory and experiment are in agreement.
		- Near-linear scaling is validated empirically.

		Presentation Bullets:
		- Graphics exponent = 1.
		- Compute exponent = 1.
		- Total exponent = 1.
		- Experimental support for near-linear scaling.
		- Results agree with complexity analysis.		
		
		
	7.5 Curvature Analysis
	
	Figure: Curvature Analysis

		Purpose:
		- Quantify departures from linear scaling.
		- Compare curvature among graphics, compute, and total execution times.
		- Assess scaling stability across the benchmark range.
		- Complement the log-log regression results.

		Supports:
		- Measured data closely follows quadratic fits.
		- Deviations from linear behavior are small.
		- Graphics execution exhibits minimal curvature.
		- Total execution exhibits minimal curvature.
		- Compute execution exhibits slightly larger curvature.
		- Curvature remains bounded in all cases.
		- Near-linear behavior is maintained throughout the benchmark range.

		Presentation Bullets:
		- Quadratic fits.
		- Curvature bounds.
		- Graphics curvature.
		- Compute curvature.
		- Total curvature.
		- Small deviations from linearity.
			Paragraph 1

		Purpose:
		- Introduce curvature analysis as a secondary scaling metric.
		- Explain why scaling exponents alone are insufficient.
		- Quantify deviations from ideal linear behavior.
		- Establish the methodology used for curvature evaluation.

		Contains:
		- Log-log regression measures overall scaling behavior.
		- Scaling exponents provide a global characterization of growth.
		- Curvature analysis provides complementary information.
		- Curvature measures deviation from ideal linear behavior.
		- Graphics execution times are analyzed.
		- Compute execution times are analyzed.
		- Total execution times are analyzed.
		- Quadratic fits are applied to the benchmark data.
		- Curvature is used to assess departures from linear scaling.

		Presentation Bullets:
		- Complement to log-log regression.
		- Measures deviation from linear behavior.
		- Graphics, compute, and total times analyzed.
		- Quadratic fitting methodology.
		- Quantifies scaling curvature.		
		


	Paragraph 2

		Purpose:
		- Present the curvature-analysis results.
		- Evaluate the magnitude of deviations from linear behavior.
		- Compare graphics, compute, and total execution-time curvature.
		- Assess the practical significance of observed curvature.

		Contains:
		- Quadratic fits and curvature bounds are presented.
		- Fitted curves closely follow the measured data.
		- Deviations from linear behavior are small.
		- Graphics execution times exhibit very small curvature.
		- Total execution times exhibit very small curvature.
		- Compute execution times exhibit somewhat larger curvature.
		- Compute curvature remains bounded.
		- No component exhibits substantial departure from linear behavior.
		- Measured behavior remains consistent with near-linear scaling.

		Presentation Bullets:
		- Quadratic fits match measured data.
		- Minor departures from linear behavior.
		- Graphics curvature is small.
		- Total curvature is small.
		- Compute curvature is slightly larger.
		- All curvature remains bounded.		
		


	Paragraph 3

		Purpose:
		- Summarize the curvature-analysis findings.
		- Relate curvature measurements to overall scaling behavior.
		- Combine curvature and regression results into a unified interpretation.
		- Reinforce the near-linear scaling conclusion.

		Contains:
		- Curvature bounds remain small.
		- Curvature is small relative to overall execution-time trends.
		- No substantial departure from linear behavior is observed.
		- Curvature results support the regression analysis.
		- Curvature and scaling exponents provide complementary evidence.
		- Both analyses indicate near-linear scaling.
		- Near-linear behavior is maintained throughout the benchmark range.
		- Experimental evidence consistently supports theoretical expectations.

		Presentation Bullets:
		- Curvature bounds are small.
		- Deviations from linearity remain minor.
		- Curvature supports regression results.
		- Multiple analyses indicate near-linear scaling.
		- Near-linear behavior maintained across the benchmark range.		

8. Discussion

8.1 Graphics-Pipeline Occupancy Generation

	Paragraph 1

		Purpose:
		- Revisit the central architectural contribution of gPCD.
		- Contrast gPCD with conventional collision-detection frameworks.
		- Explain how rendering and occupancy generation are organized differently.
		- Establish the foundation for the discussion that follows.

		Contains:
		- Traditional frameworks separate collision detection and rendering.
		- Rendering typically consumes simulation state.
		- Simulation data is passed to the graphics pipeline for visualization.
		- Occupancy generation is traditionally a separate operation.
		- gPCD integrates occupancy generation into the graphics pipeline.
		- Occupancy generation and rendering occur concurrently.
		- The graphics pipeline participates directly in collision detection.
		- The framework adopts a different architectural organization than conventional approaches.

		Presentation Bullets:
		- Traditional approach:
			• Simulation
			• Collision detection
			• Rendering

		- gPCD approach:
			• Occupancy generation in graphics pipeline
			• Rendering and occupancy construction proceed concurrently

		- Architectural integration of rendering and collision detection.



	Paragraph 2

		Purpose:
		- Clarify the nature of the gPCD contribution.
		- Distinguish architectural novelty from algorithmic novelty.
		- Explain how graphics hardware is repurposed within the framework.
		- Reinforce the execution-model contribution of the work.

		Contains:
		- The primary contribution is architectural.
		- The contribution is not a new collision-detection algorithm.
		- The contribution is not a new spatial data structure.
		- gPCD introduces a new execution model.
		- Graphics-pipeline operations participate directly in broad-phase collision detection.
		- Occupancy generation is integrated into graphics execution.
		- Graphics hardware performs work beyond visualization.
		- Rendering resources contribute directly to collision detection.
		- The framework repurposes resources traditionally reserved for rendering.

		Presentation Bullets:
		- Architectural contribution.
		- Not a new data structure.
		- Not a new broad-phase methodology.
		- New execution model.
		- Graphics pipeline contributes to collision detection.
		- Rendering resources perform occupancy generation.




	Paragraph 3

		Purpose:
		- Discuss the synchronization characteristics of gPCD.
		- Explain how the framework reduces coordination requirements.
		- Relate the architecture to massively parallel execution.
		- Highlight an implementation advantage of the particle-centric design.

		Contains:
		- Synchronization requirements are limited.
		- Most synchronization occurs during occupancy generation.
		- Atomic operations are used only for occupancy-list insertion.
		- Shared-cell slot reservation requires thread-safe insertion.
		- Collision discovery is particle-centric.
		- No global collision-pair structure is constructed.
		- Candidate ownership remains local to the source particle.
		- Narrow-phase evaluation maps naturally onto GPU parallelism.
		- The architecture minimizes global coordination requirements.
		- The framework is well suited to massively parallel execution.

		Presentation Bullets:
		- Synchronization localized to occupancy generation.
		- Atomics only for cell-list insertion.
		- No global collision-pair structure.
		- Particle-centric collision discovery.
		- Natural GPU parallelism.
		- Reduced synchronization overhead.



	Paragraph 4

		Purpose:
		- Relate experimental performance results to theoretical complexity.
		- Synthesize evidence from multiple performance analyses.
		- Evaluate whether measured behavior matches theoretical expectations.
		- Discuss the practical implications of the scaling results.

		Contains:
		- Throughput results are consistent with theoretical expectations.
		- Time-per-particle results are consistent with theoretical expectations.
		- Log-log regression results are consistent with theoretical expectations.
		- Curvature results are consistent with theoretical expectations.
		- Measured behavior agrees with O(N) occupancy-generation complexity.
		- Implementation-specific effects influence absolute execution times.
		- Hardware and implementation details affect constant factors.
		- No dominant super-linear behavior was observed.
		- Measured scaling remained stable across the benchmark range.
		- Experimental results support the complexity analysis.

		Presentation Bullets:
		- Throughput supports theory.
		- Time-per-particle supports theory.
		- Regression supports theory.
		- Curvature supports theory.
		- Consistent with O(N) occupancy generation.
		- No dominant super-linear behavior observed.
		- Constant factors affect absolute performance.

	8.2 Scalability of the gPCD Framework

	Paragraph 1

		Purpose:
		- Summarize the scalability evidence presented in the Results section.
		- Synthesize multiple performance analyses into a unified conclusion.
		- Relate measured behavior to the theoretical complexity analysis.
		- Establish the basis for the scalability discussion that follows.

		Contains:
		- The framework exhibits near-linear scaling.
		- Scaling behavior is maintained throughout the evaluated benchmark range.
		- Throughput results support the scalability conclusion.
		- Time-per-particle results support the scalability conclusion.
		- Log-log regression results support the scalability conclusion.
		- Curvature results support the scalability conclusion.
		- Multiple independent analyses are consistent.
		- Measured behavior agrees with the complexity analysis.
		- Experimental evidence supports the scalability claims.

		Presentation Bullets:
		- Near-linear scaling observed.
		- Throughput supports scalability.
		- Time-per-particle supports scalability.
		- Regression supports scalability.
		- Curvature supports scalability.
		- Results agree with complexity analysis.



	Paragraph 2

		Purpose:
		- Interpret the measured scaling and curvature results.
		- Examine scalability at the level of individual execution stages.
		- Identify potential sources of nonlinear growth.
		- Explain why near-linear behavior is maintained.

		Contains:
		- Graphics scaling exponent remains near unity.
		- Compute scaling exponent remains near unity.
		- Total execution-time exponent remains near unity.
		- Curvature bounds remain small.
		- Graphics execution exhibits stable scaling behavior.
		- Compute execution exhibits stable scaling behavior.
		- No stage exhibits strong nonlinear growth.
		- No single stage dominates overall scaling behavior.
		- Performance growth remains balanced across framework components.
		- Near-linear overall scaling emerges from near-linear behavior within individual stages.

		Presentation Bullets:
		- Graphics exponent = 1.
		- Compute exponent = 1.
		- Total exponent = 1.
		- Small curvature bounds.
		- No dominant nonlinear stage.
		- Balanced scaling across framework components.

	Paragraph 3

		Purpose:
		- Distinguish throughput saturation from scaling degradation.
		- Explain the behavior observed at large particle counts.
		- Address a potential misinterpretation of the throughput results.
		- Clarify the relationship between utilization limits and scalability.

		Contains:
		- Throughput eventually enters a throughput-limited regime.
		- Throughput behavior changes at large particle counts.
		- The throughput limitation is associated with resource utilization.
		- GPU resources approach sustained operating capacity.
		- Throughput saturation does not imply a loss of scalability.
		- Underlying scaling behavior remains stable.
		- The framework continues to exhibit near-linear characteristics.
		- The observed behavior reflects hardware utilization limits rather than algorithmic degradation.
		- Large-scale performance remains consistent with the complexity analysis.

		Presentation Bullets:
		- Throughput-limited regime observed.
		- High resource utilization.
		- Throughput saturation is expected.
		- Not evidence of scaling breakdown.
		- Underlying scaling behavior remains stable.
		- Consistent with complexity analysis.

	8.3 Memory Locality Considerations

	Paragraph 1

		Purpose:
		- Interpret the memory-locality benchmark results.
		- Distinguish complexity from practical performance.
		- Explain why memory organization influences execution time.
		- Relate locality effects to framework performance.

		Contains:
		- Memory locality influences measured performance.
		- Memory organization affects execution time.
		- Locality effects occur independently of asymptotic complexity.
		- Sequential and random layouts have the same theoretical complexity.
		- Sequential ordering improves memory locality.
		- Improved locality reduces execution time.
		- The locality benchmark isolates memory-access effects.
		- Performance differences arise from implementation characteristics rather than algorithmic complexity.

		Presentation Bullets:
		- Memory locality matters.
		- Same asymptotic complexity.
		- Different execution times.
		- Sequential ordering improves locality.
		- Improved locality reduces runtime.
		- Complexity and performance are distinct considerations.



	Paragraph 2

		Purpose:
		- Interpret the locality benchmark trends.
		- Relate observed behavior to scaling characteristics.
		- Distinguish locality effects from asymptotic behavior.
		- Identify the primary impact of memory locality.

		Contains:
		- Sequential and random configurations exhibit similar trends.
		- The execution-time difference remains approximately proportional across the benchmark range.
		- The observed behavior suggests locality affects constant factors.
		- Memory-access efficiency is the primary source of the performance difference.
		- Locality does not substantially alter scaling behavior.
		- The underlying growth characteristics remain unchanged.
		- Locality influences implementation efficiency.
		- Asymptotic framework behavior remains unaffected.
		- Performance differences arise from practical execution effects rather than complexity differences.

		Presentation Bullets:
		- Approximately parallel trends observed.
		- Locality affects memory-access cost.
		- Constant factors are influenced.
		- Scaling behavior remains unchanged.
		- Implementation efficiency is affected.
		- Asymptotic behavior is preserved.


	8.4 Implications for Large-Scale Particle Simulation
	
	Figure: Particle Density Comparison

		Purpose:
		- Visualize the effect of increasing particle count.
		- Demonstrate increasing spatial resolution.
		- Compare sparse and dense particle representations.
		- Illustrate the practical value of scalability.

		Supports:
		- Low particle counts exhibit visible discreteness.
		- High particle counts approach continuum-like appearance.
		- Particle density increases spatial fidelity.
		- Large particle counts enable more detailed physical representations.
		- Scalability directly translates into simulation capability.

		Presentation Bullets:
		- 5,832 particles.
		- 2,985,984 particles.
		- Resolution increase.
		- Continuum-like appearance.
		- Retained discrete formulation.

	Paragraph 1

		Purpose:
		- Relate scalability to physical modeling capability.
		- Explain the practical significance of increasing particle count.
		- Connect framework performance to simulation fidelity.
		- Motivate the broader applicability of gPCD.

		Contains:
		- Increasing particle count increases spatial resolution.
		- Low particle counts reveal the discrete nature of the representation.
		- High particle counts approach visually continuous behavior.
		- The formulation remains fundamentally discrete.
		- Increased particle density improves physical representation.
		- Higher particle counts enable investigation of more complex phenomena.
		- Large-scale particle systems can model a broad range of physical processes.
		- Macroscopic behavior emerges from interacting discrete particles.
		- The demonstrated scalability expands the range of feasible particle-based simulations.

		Presentation Bullets:
		- Higher particle counts increase spatial resolution.
		- Discrete systems approach continuum-like behavior.
		- Representation remains particle based.
		- Increased simulation fidelity.
		- Broader physical applicability.



	Paragraph 2

		Purpose:
		- Discuss the practical usability of high-density particle configurations.
		- Relate particle density to interactive simulation performance.
		- Explain the significance of maintaining real-time execution.
		- Connect collision-detection performance to broader simulation capability.

		Contains:
		- High-density particle configurations remain interactive.
		- Approximately three million particles sustain about 50 FPS.
		- Large particle counts do not preclude interactive visualization.
		- Computational resources remain available beyond collision detection.
		- Particle dynamics can be evaluated concurrently.
		- Force calculations remain feasible.
		- Additional simulation operations remain feasible.
		- High particle density and practical simulation performance can coexist.
		- The framework supports both large-scale representation and interactive execution.

		Presentation Bullets:
		- ~3 million particles.
		- ~50 FPS.
		- Interactive performance maintained.
		- Resources remain available for simulation.
		- Dynamics and force evaluation remain feasible.
		- Large-scale simulations remain practical.

	Paragraph 3

		Purpose:
		- Interpret the demonstrated particle counts in practical terms.
		- Relate benchmark results to simulation scale.
		- Position the framework relative to traditional HPC workloads.
		- Highlight the significance of achieving large-scale simulation on a single CPU-GPU system.

		Contains:
		- The visualizations provide scale intuition.
		- Demonstrated particle counts represent practically significant system sizes.
		- Large particle counts can be achieved on a single CPU-GPU platform.
		- The framework supports simulation scales traditionally associated with HPC resources.
		- High particle counts are no longer restricted to large computing clusters.
		- Large-scale particle simulation becomes accessible on workstation-class hardware.
		- The demonstrated scalability extends the practical applicability of particle-based methods.
		- The framework lowers the hardware requirements associated with large particle systems.

		Presentation Bullets:
		- Visual scale comparison.
		- Demonstrated multi-million-particle systems.
		- Single CPU-GPU execution.
		- HPC-scale particle counts.
		- Workstation-class accessibility.
		- Expanded simulation capability.

9. Conclusion and Future Work

	Paragraph 1

		Purpose:
		- Summarize the primary contribution of the work.
		- Restate the architectural innovation introduced by gPCD.
		- Contrast gPCD with conventional collision-detection frameworks.
		- Reiterate the role of graphics-pipeline occupancy generation.

		Contains:
		- gPCD is presented as a graphics-based parallel collision-detection framework.
		- Broad-phase occupancy generation is integrated into the graphics pipeline.
		- Traditional frameworks treat rendering as a consumer of simulation state.
		- gPCD assigns an active collision-detection role to the graphics pipeline.
		- Occupancy construction and rendering proceed concurrently.
		- The framework generates the potentially colliding set during graphics execution.
		- Candidate generation supports subsequent narrow-phase evaluation.
		- The primary contribution is architectural organization rather than a new data structure.

		Presentation Bullets:
		- Graphics-based Parallel Collision Detection (gPCD).
		- Graphics-pipeline occupancy generation.
		- Concurrent rendering and occupancy construction.
		- PCS generation during graphics execution.
		- Architectural contribution.


	Paragraph 2

		Purpose:
		- Summarize the analytical and experimental findings.
		- Connect theoretical predictions to measured behavior.
		- Reinforce the scalability conclusions.
		- Consolidate the major results of the study.

		Contains:
		- Complexity analysis was developed for all major framework components.
		- Occupancy generation complexity was analyzed.
		- Candidate construction complexity was analyzed.
		- Narrow-phase complexity was analyzed.
		- Storage requirements were analyzed.
		- Experimental evaluation demonstrated high throughput.
		- Experimental evaluation demonstrated near-linear scaling.
		- Graphics execution exhibited near-linear behavior.
		- Compute execution exhibited near-linear behavior.
		- Total execution exhibited near-linear behavior.
		- Log-log regression supported the scaling conclusions.
		- Curvature analysis supported the scaling conclusions.
		- Measured behavior was consistent with theoretical expectations.
		- Locality studies identified memory-access effects.
		- Memory locality influenced implementation efficiency.
		- Memory locality did not alter asymptotic behavior.

		Presentation Bullets:
		- Complexity analysis completed.
		- High throughput demonstrated.
		- Near-linear scaling demonstrated.
		- Regression supports theory.
		- Curvature supports theory.
		- Locality affects efficiency.
		- Asymptotic behavior preserved.

	Paragraph 3

		Purpose:
		- State the broader implications of the results.
		- Evaluate the viability of the proposed architecture.
		- Connect measured performance to practical capability.
		- Position the framework relative to traditional HPC expectations.

		Contains:
		- Graphics-pipeline occupancy generation is a viable execution architecture.
		- The proposed organization successfully supports large-scale collision detection.
		- Experimental results demonstrate practical scalability.
		- Large particle counts can be achieved on a single CPU-GPU system.
		- Multi-million-particle simulations are feasible on workstation-class hardware.
		- The demonstrated scales have traditionally been associated with HPC environments.
		- The framework expands the accessibility of large-scale particle simulation.
		- The architectural approach provides a practical alternative to conventional execution models.

		Presentation Bullets:
		- Graphics-pipeline occupancy generation is viable.
		- Large-scale collision detection demonstrated.
		- Single CPU-GPU execution.
		- Multi-million-particle capability.
		- Traditionally HPC-associated scales.
		- Increased accessibility of large-scale simulation.

	Paragraph 4

		Purpose:
		- Identify future research directions.
		- Discuss planned extensions of the framework.
		- Connect collision detection to broader simulation applications.
		- Position the framework within future thermo-fluid and multiphysics research.

		Contains:
		- Future work includes larger GPU configurations.
		- Future work includes multi-GPU execution.
		- Future work includes distributed computing environments.
		- Current work focuses primarily on collision detection.
		- The framework has already been applied to particle-based thermo-fluid simulations.
		- Preliminary thermo-fluid studies have produced promising results.
		- Novel physical observations have emerged from preliminary investigations.
		- Direct simulation remains important for scientific discovery.
		- Machine-learning methods depend on representative training data.
		- Simulation-generated data can support data-driven models.
		- Future work will investigate thermo-fluid applications in greater depth.
		- Future work will investigate multiphysics applications.
		- Future work will evaluate graphics-pipeline occupancy generation in broader simulation domains.

		Presentation Bullets:
		- Larger GPU configurations.
		- Multi-GPU execution.
		- Distributed computing.
		- Thermo-fluid simulation.
		- Multiphysics simulation.
		- Novel physical observations.
		- Simulation-generated training data.
		- AI-assisted predictive modeling.




   \end{verbatim}

